% Generated draft: requires mathematical and source review.
\documentclass[11pt]{article}
\usepackage[margin=25mm]{geometry}
\usepackage{fontspec}
\setmainfont{Latin Modern Roman}
\usepackage{amsmath,amssymb,mathtools,unicode-math}
\setmathfont{Latin Modern Math}
\usepackage{enumitem,needspace,xurl,hyperref}
\setlength{\parindent}{0pt}
\setlength{\parskip}{5pt}
\setlength{\emergencystretch}{3em}
\newcommand{\N}{\mathbb{N}}
\newcommand{\Z}{\mathbb{Z}}
\newcommand{\Q}{\mathbb{Q}}
\newcommand{\R}{\mathbb{R}}
\newcommand{\C}{\mathbb{C}}
\newcommand{\F}{\mathbb{F}}
\newcommand{\A}{\mathbb{A}}
\newcommand{\PP}{\mathbb{P}}
\newcommand{\E}{\mathbb{E}}
\newcommand{\eps}{\varepsilon}
\newcommand{\dd}{\,\mathrm{d}}
\begin{document}
\section*{UnsolvedMath ID: 700876}
\textbf{Source problem code:} catalogue entry
\subsection{Definitions and mathematical statement}
Let \(G\) be a finite group and let \(A\leq G\) be cyclic of order \(4\). For \(g\in G\), write \(A^g=g^{-1}Ag\). Say that \(A\) is a \(TI\)-subgroup if
\[
\forall g\in G,\qquad A\cap A^g\in\{\{1\},A\}.
\]
In particular, no assumption that \(A\) is nonnormal is imposed. A component of \(G\) is a subnormal quasisimple subgroup \(L\leq G\): thus there is a chain \(L=L_0\trianglelefteq L_1\trianglelefteq\cdots\trianglelefteq L_r=G\), and \(L=[L,L]\) with \(L/Z(L)\) a nonabelian simple group. Here \(Z(L)\) denotes the center of \(L\). Write
\[
C_G(L)=\{g\in G: g\ell=\ell g\text{ for every }\ell\in L\}.
\]
The source notation \(C(L)\) is interpreted as \(C_G(L)\); the ambient centralizer is not explicitly defined in the supplied record. The product \(LC_G(L)\) is a subgroup of \(G\).

\begin{enumerate}
\item[(a)] Is it true that, for every such pair \((G,A)\) and every component \(L\) of \(G\),
\[
A\cap LC_G(L)=\{1\}\quad\Longrightarrow\quad A\leq C_G(L)?
\]
No additional hypothesis that \(A\) normalizes \(L\) is supplied.
\item[(b)] Assuming that \(A\leq C_G(L)\) for every component \(L\) of \(G\), determine the structure of the normal closure
\[
\langle A^G\rangle=\bigl\langle A^g:g\in G\bigr\rangle.
\]
The assumption includes groups with no components, for which it is vacuous. The source does not specify the desired form of the structural description or classification.
\end{enumerate}
\subsection{Short English statement}
For a cyclic \(TI\)-subgroup \(A\) of order \(4\) in a finite group \(G\), ask whether trivial intersection with \(LC_G(L)\) forces \(A\) to centralize a component \(L\). Also determine the structure of the normal closure of \(A\) when \(A\) centralizes every component of \(G\).
\subsection{Sources}
\begin{itemize}
\item Local UnsolvedMath registry, supplied complete source record: permanent ID \(700876\), exact source problem code \texttt{catalogue entry}, title \emph{Kourovka Notebook Problem 12.57}. The registry labels the problem \texttt{open}; this is a supplied status, not an independently verified mathematical or literature assessment.
\item Attributed in the supplied record to A. A. Makhn\"ev, proposed in \(1992\), as Problem 12.57 of \emph{The Kourovka Notebook (Unsolved Problems in Group Theory)}, 12th issue (1992). The record also identifies the 21st issue edition (2026), \texttt{arXiv:1401.0300v48}. These bibliographic details are supplied attribution and have not been independently verified.
\item The supplied collection description names E. I. Khukhro and V. D. Mazurov as editors of \emph{Kourovka Notebook (Issues 1--21)}, covering 1965--2026. No separate literature-verification date or additional literature-source links are supplied.
\end{itemize}
\end{document}
